\documentclass[aps,prd,twocolumn,superscriptaddress,nofootinbib]{revtex4-2}

\usepackage{graphicx}
\usepackage{amsmath}
\usepackage{amssymb}
\usepackage{hyperref}

\begin{document}

\title{Spectral analysis of systematic residuals in Casimir force measurements via atomic force microscopy}

\author{Bruno Lu\c{c}ardo Olivaes}
\email{bruno@example.com} % Update with your actual email
\affiliation{Independent Researcher} % Update with appropriate affiliation

\date{\today}

\begin{abstract}
Precision measurements of the Casimir force between macroscopic bodies have historically been challenged by systematic residuals, particularly at short separation distances. Standard theoretical models attribute these deviations primarily to surface roughness and finite-temperature corrections. In this paper, we present a novel methodology for the spectral analysis of raw atomic force microscopy (AFM) data to investigate these short-distance anomalies. By applying Fast Fourier Transform (FFT) protocols to existing historical datasets, we test the hypothesis that the observed residuals contain periodic components which are not adequately explained by conventional surface models. We aim to establish whether these periodic signals represent uncharacterized instrumental artifacts or point towards novel, short-range interactions. Our results provide a refined protocol for analyzing AFM force-distance curves and propose new constraints on theoretical models of the Casimir effect at sub-micron scales.
\end{abstract}

\maketitle

\section{Introduction}
% Introduction goes here...

\section{Methodology}
\subsection{Atomic Force Microscopy Data Acquisition}
% Details on data from Lamoreaux, Decca, etc.

\subsection{FFT Protocol for Force-Distance Curves}
% Explanation of the spectral analysis method.

\section{Results}
% To be populated with FFT analysis results.

\section{Discussion}
% Comparison with standard roughness/temperature models.

\section{Conclusion}
% Concluding remarks.

\begin{acknowledgments}
We acknowledge the use of historical datasets provided by...
\end{acknowledgments}

\bibliography{references}

\end{document}
